\chapter{Обход графа: один движок, три политики}
\label{ch:walker}

\section{Случайное блуждание как универсальный механизм}

Почти всё, что Рин делает со своей памятью, — это \emph{обход} графа роем случайных
блуждающих агентов, распространяющих активацию. Retrieval, внутренняя жизнь и сон —
не три разных движка, а \emph{три политики} одного движка, различающиеся смещением
перехода, вероятностью телепорта и набором операций.

Общее правило перехода — нормированное произведение факторов:
\begin{equation}
P(A\to B)\;\propto\;
\underbrace{w_{AB}}_{\text{сила связи}}\;\cdot\;
\underbrace{\text{rel}(B\mid q)}_{\text{релевантность запросу}}\;\cdot\;
\underbrace{\text{mc}(B\mid \mood)}_{\text{конгруэнтность настроению}}\;\cdot\;
\underbrace{\text{unc}(B)^{\beta_u}}_{\text{тяга к неопределённости}},
\label{eq:trans}
\end{equation}
где $\text{mc}$ — из \eqref{eq:moodcong}, а показатель $\beta_u$ включает/выключает
любопытство.

\begin{intuition}[Что такое «температура»]
«Температура» — одна ручка: насколько выбор следующего шага случаен против жадного.
Формально переход берётся softmax-ом от логарифма правой части \eqref{eq:trans},
делённого на температуру $T$: при $T\to 0$ — жадный шаг по сильнейшему ребру
(эксплуатация), при большом $T$ — почти равновероятный выбор соседа (исследование).
Термин из имитации отжига; мистики нет.
\end{intuition}

\section{Три политики как параметризация}

\begin{center}
\begin{tabular}{p{0.24\textwidth} p{0.10\textwidth} p{0.10\textwidth} p{0.10\textwidth} p{0.30\textwidth}}
\toprule
\textbf{Политика} & $\alpha$ (телепорт) & $\beta_u$ & $T$ & \textbf{Что запускает} \\
\midrule
retrieval (цель) & $0$ & $0$ & низкая & извлечение под запрос \\
внутренняя жизнь (любопытство) & мал. & $>0$ & средняя & inspect, fusion у неопределённых \\
сон (покрытие) & большой & $0$ & высокая & fusion/inspect у случайных \\
\bottomrule
\end{tabular}
\end{center}

\paragraph{Retrieval — целенаправленная.} Стартует от заякоренных запросом узлов,
собирает \emph{окрестность} активированного содержания. Подчеркнём (глава~\ref{ch:personality}):
обход собирает \emph{ингредиенты}, а не готовую реакцию.

\paragraph{Внутренняя жизнь — по неопределённости.} Смещена к узлам с высоким
\texttt{need\_to\_discover}, низкой уверенностью, новизной ($\beta_u>0$).
«Любопытство» — целенаправленный поиск неопределённости, не случайность.

\paragraph{Сон — телепортирующий обходчик.} Политика \emph{стохастического
покрытия}: с вероятностью $\alpha$ телепорт в случайный узел, иначе — шаг по
\eqref{eq:trans}. Это случайное блуждание с рестартом (персонализированный
PageRank). Его стационарное распределение удовлетворяет
\begin{equation}
\pi \;=\; (1-\alpha)\,P^\top \pi \;+\; \alpha\, \mathbf{r},
\label{eq:ppr}
\end{equation}
где $\mathbf{r}$ — распределение рестарта. Телепорт делает блуждание эргодическим —
это ключевое свойство для эпистемической гигиены (глава~\ref{ch:topology}).

\begin{intuition}[Три способа гулять по городу]
Retrieval — идти по адресу (есть цель, выбираешь осмысленный путь). Внутренняя
жизнь — идти туда, где интересно и непонятно (тянет к незнакомым переулкам). Сон —
время от времени садиться в случайное такси и ехать в случайный район: не потому,
что там дело, а чтобы не застревать в одном квартале. Один пешеход, три правила
выбора улицы.
\end{intuition}

\section{Фоновая разведка: один процесс, а не «день/ночь»}
\label{sec:background}

Разделение на «дневное любопытство» и «ночной обходчик» — удобная, но неточная
рамка. Точнее: есть \emph{один всегда включённый} фоновый разведчик, а «день» и
«ночь» — два конца \emph{континуума по свободным ресурсам}, а не переключатель.

Формализуем «мощности, внимание, фокус». Пусть $\rho(t)\in[0,1]$ —
\emph{foreground-нагрузка}: насколько текущий момент требует активного ответа
($\rho\approx 1$ в напряжённом диалоге, $\rho\approx 0$ в простое). Общий бюджет
$B_{\text{total}}$ делится с приоритетом переднего плана; фоновой разведке достаётся
\begin{equation}
B_{\text{bg}}=B_{\text{total}}\,(1-\rho).
\end{equation}
\emph{Фокус} — это активное множество $A_{\text{focus}}$ (узлы высокой активации
$a_v$). Параметры фонового роя — функции $B_{\text{bg}}$: число walker'ов
$\propto B_{\text{bg}}$; доля телепортов $\alpha$ и «широта» растут при падении
$\rho$; стартовое распределение интерполирует
\[
r=(1-\beta)\cdot\mathrm{anchor}(A_{\text{focus}})+\beta\cdot\mathrm{uniform},
\qquad \beta\uparrow \text{ при } \rho\downarrow.
\]
Тогда прежние «любопытство» и «сон» — это два режима одного процесса:
\emph{высокий} $\rho$ (мало бюджета) $\Rightarrow$ мало walker'ов, $\alpha\approx 0$,
старт от фокуса — локальная разведка вокруг текущей темы; \emph{низкий} $\rho$ (много
бюджета) $\Rightarrow$ много walker'ов, большой $\alpha$, старт почти равномерный —
широкий телепорт-скраб (глава~\ref{ch:sleep}).

\begin{intuition}[Мысль на перекуре]
Именно поэтому яркая ассоциация может прийти \emph{днём}, «на ровном месте»: фоновый
разведчик работает всегда, и даже на скудном дневном бюджете иногда набредает на
далёкую связь. Развёрнутый пример такой дневной находки — в главе~\ref{ch:fusion}
(House\,/\,Cast Away).
\end{intuition}

\begin{openq}[Как измерять foreground-нагрузку $\rho$]
$\rho$ управляет тем, насколько агрессивно разведчик забирает ресурсы, но как её
\emph{мерить} в рантайме — не зафиксировано. Кандидаты: наличие активного
пользовательского хода (есть незакрытая реплика $\Rightarrow$ $\rho$ высока); длина
очереди foreground-задач; «важность» текущего события (по appraisal,
глава~\ref{ch:personality}); загрузка вычислителя. Вероятно, комбинация, но веса и
динамика (как быстро $\rho$ спадает в простое) пока открыты.
\end{openq}

\section{Единство и следствие}

Один субстрат, один блуждающий агент, три параметризации: смещение перехода,
вероятность телепорта, набор операций. Реализуется одна машина обхода с параметрами
(псевдокод — приложение~\ref{ch:algorithms}). Родственные механизмы: spreading
activation (Collins–Loftus) как модель ассоциативного припоминания и
персонализированный PageRank / RWR как его вычислительная форма; в волте RIN — под
именами SYNAPSE и GAAMA.
